\begin{trapbox}
\begin{enumerate}
    \item \textbf{把最大张角误写成 $2\arcsin(b/a)$。}
    短轴端点的半角直角三角形中，对边是 $c$，斜边是 $a$，因此
    \[
    \sin\frac{\theta_{\max}}2=\frac ca=e,
    \]
    而不是 $b/a$。$b/a$ 对应的是该直角三角形的另一个锐角的正弦。

    \item \textbf{误以为“任意点都取等”。}
    $\sin(\theta/2)\le e$ 对所有椭圆上的点成立，但只有
    $P=(0,\pm b)$ 取等号；长轴端点 $P=(\pm a,0)$ 反而有 $\theta=0$。

    \item \textbf{忽略三角函数的单调区间。}
    因为椭圆的两个焦点不同，且 $P$ 在椭圆边界，
    $0\le\theta<\pi$，所以 $0\le\theta/2<\pi/2$。
    $\sin(\theta/2)$ 在这个范围内严格随 $\theta$ 增大，
    才能把半角不等式转成 $\theta\le2\arcsin e$。

    \item \textbf{基本不等式的方向弄反。}
    和固定时 $r_1r_2\le a^2$，因此
    \[
    \frac{b^2}{r_1r_2}\ge\frac{b^2}{a^2},
    \quad
    1-\frac{b^2}{r_1r_2}\le1-\frac{b^2}{a^2}.
    \]
    分母越大，分数越小；前面还有减号，方向需要特别检查。

    \item \textbf{“存在直角”与“存在钝角”的边界混淆。}
    \[
    \begin{array}{ll}
    \text{存在直角：}&e\ge\sqrt2/2,\\
    \text{存在钝角：}&e>\sqrt2/2,\\
    \text{所有点的焦点张角都不大于直角：}&e\le\sqrt2/2.
    \end{array}
    \]
    注意：存在性的条件不能随意改写成“每一个点都满足”。

    \item \textbf{漏掉椭圆离心率本身的范围。}
    标准非圆椭圆要求 $0<e<1$。例如 $e\ge\sqrt3/2$
    的最终范围应是 $[\sqrt3/2,1)$，右端的 $1$ 不能取。

    \item \textbf{直接照搬到双曲线。}
    原文件写“对双曲线有 $\cos(\theta/2)\le1/e$”，
    这\textbf{不是普遍正确的不等式}，应删除。
    例如双曲线右支上 $P$ 沿曲线向无穷远移动时，
    两条指向焦点的视线越来越接近，$\theta\to0$，于是
    $\cos(\theta/2)\to1$；然而双曲线 $e>1$，所以 $1/e<1$，
    上述所谓不等式会被违反。
    椭圆的关键前提是 $PF_1+PF_2=2a$；
    双曲线使用的是焦半径\emph{差}的绝对值等于 $2a$，不能套同一证明。
\end{enumerate}
\end{trapbox}
